%%%%%%%%%%%%%%%%%%%%%%%%%%%%%%%%%%%%%%%%%
% Memo
% LaTeX Template
% Version 1.0 (30/12/13)
%
% This template has been downloaded from:
% http://www.LaTeXTemplates.com
%
% Original author:
% Rob Oakes (http://www.oak-tree.us) with modifications by:
% Vel (vel@latextemplates.com)
%
% License:
% CC BY-NC-SA 3.0 (http://creativecommons.org/licenses/by-nc-sa/3.0/)
%
%%%%%%%%%%%%%%%%%%%%%%%%%%%%%%%%%%%%%%%%%

\documentclass[letterpaper,11pt]{texMemo} % Set the paper size (letterpaper, a4paper, etc) and font size (10pt, 11pt or 12pt)

\usepackage{fancyhdr}
\usepackage{fancybox}
\usepackage{longtable}
\usepackage{amsmath}
\usepackage{xcolor}
%----------------------------------------------------------------------------------------
%	MEMO INFORMATION
%----------------------------------------------------------------------------------------



\memoto{Luis Andr\'es Valido Fajardo. luis.valido@umcc.cu (53694742)} % Recipient(s)

\memofrom{Josval Díaz Blanco} % Sender(s)

\memosubject{Guía de Aprendizaje para Concursantes ICPC y IOI: Búsqueda Binaria } % Memo subject

\memodate{\today} % Date, set to \today for automatically printing todays date

\logo{\includegraphics[scale=0.5]{img/icpc}} % Institution logo at the top right of the memo, comment out this line for no logo

%----------------------------------------------------------------------------------------

%\titleguide{Detección de ciclos en un grafo}%cycle_graph
%\titleguide{Camino y ciclo de Hamilton}%hamilton_tour
%\titleguide{Encontrarse en el medio ( \emph{Meet in the Middle} )}%meet_in_middle
%\titleguide{Números Catalanes}%number_catalan
%\titleguide{Introducción al ajedrez}%introduction_chess 
%\titleguide{Operaciones con matrices}%matrix_operation 
%\titleguide{Coeficientes binomiales}%binomial_coefficients
%\titleguide{Funciones de distancias}%_distance_functions
%\titleguide{Tabla dispersa (\emph{Sparse Table})}%sparse_table 
%\titleguide{Algoritmo-Z (\emph{Z-Algorithm})}%z_algorithm 
%\titleguide{Función de prefijo. Algoritmo de Knuth-Morris-Pratt (\emph{KMP})}%kmp
%\titleguide{Descomposición raíz cuadrada (\emph{Sqrt Decomposition})}%sqrt_decomposition
%\titleguide{Algoritmo de Mo (\emph{Mo's algorithm})}%mo_algorithm 
%\titleguide{Probabilidades}%probability
%\titleguide{Funciones de String en C++}%f_string_cpp
%\titleguide{Arreglos de sufijos (\emph{Suffix Array})}%suffix_array FALTA
%\titleguide{Permutaciones}%permutation FALTA
%\titleguide{Sucesiones númericas}%successions FALTA
%\titleguide{Arbol de raíz cuadrada (\emph{SQRT Tree})}%sqrt_tree FALTA
%\titleguide{Cadenas de carácteres, terminología}%_string_terminology FALTA
\titleguide{Búsqueda ternaria}%_ternary_search
\begin{document}

%\AddToShipoutPicture{\BackgroundPic}
\maketitle % Print the memo header information
%\tableofcontents
\pagestyle{plain}
\pagebreak

\pagestyle{fancy}
\fancyhead[LO,CE]{ }
\fancyhead[RO,CE]{\includegraphics[scale=0.1]{img/icpc}}
\fancyfoot[LO,CE]{\textbf{Autor:} Luis Andrés Valido Fajardo \\ \textbf{Email:} luis.valido1989@gmail.com \\ \textbf{Teléfono:} 53694742}
\fancyfoot[RO,CE]{\emph{Existen 10 tipos de personas Las que \\saben binario y LAS QUE NO}}
\fancypagestyle{plain}{\pagestyle{fancy}}



%\lhead{ }
%\rhead{  }

%\fancyfoot[L]{}
%\fancyfoot[R]{\textbf{Autor:} Luis Andrés Valido Fajardo \\ \textbf{Email:} luis.valido@umcc.cu}
%----------------------------------------------------------------------------------------
%	MEMO CONTENT
%----------------------------------------------------------------------------------------


\section{Introducción}
Se nos da una función $f(x)$ que es \textbf{unimodal} en un intervalo $[l,r]$. En esta guía veremos como hallar el máximo o mínimo valor según sea la situación dentro del intervalo $[l,r]$
\section{Conocimientos previos}
\subsection{Función unimodal}
Una función unimodal es una función matemática que tiene un único máximo o mínimo absoluto en un intervalo determinado. Esto significa que la función aumenta o disminuye hasta alcanzar un punto máximo o mínimo y luego vuelve a cambiar de dirección. En otras palabras, la función tiene una forma de \emph{campana} con un solo pico o valle. Por función unimodal, nos referimos a uno de dos comportamientos de la función:

\begin{enumerate}
	\item La función primero aumenta estrictamente, alcanza un máximo (en un solo punto o en un intervalo) y luego disminuye estrictamente.
	\item La función primero disminuye estrictamente, alcanza un mínimo y luego aumenta estrictamente.
\end{enumerate}

\subsection{Teorema de Maestro}
En el análisis de algoritmos , el teorema maestro de las recurrencias de divide y vencerás proporciona un análisis asintótico para muchas relaciones de recurrencia que ocurren en el análisis de los algoritmos de divide y vencerás . El enfoque fue presentado por primera vez por Jon Bentley , Dorothea Blostein (de soltera Haken) y James B. Saxe en 1980, donde se describió como un método unificador para resolver tales recurrencias. El nombre \emph{teorema maestro} fue popularizado por el libro de texto de algoritmos ampliamente utilizado Introducción a los algoritmos de Cormen , Leiserson , Rivest y Stein .
\section{Desarrollo}
Considere 2 puntos cualesquiera $m_1$, y $m_2$ en este intervalo: $l<m_1 <m_2 <r$. Evaluamos la función en $m_1$ y $m_2$, es decir, encontrar los valores de $f(m_1)$ y $f(m_2)$. Ahora, tenemos una de tres opciones:

\begin{enumerate}
	\item $f(m_1) < f(m_2)$ El máximo deseado no se puede ubicar en el lado izquierdo de 
	$m_1$, es decir, en el intervalo $[l,m_1]$, ya que ambos puntos $m_1$ y $m_2$ o solo 
	$m_1$ pertenecen al área donde la función aumenta. En cualquier caso, esto significa 
	que tenemos que buscar el máximo en el intervalo $[m_1,r]$.
	\item $f(m_1) > f(m_2)$ Esta situación es simétrica a la anterior: el máximo no puede ubicarse en el lado derecho de $m_2$, es decir, en el intervalo $[m_2,r]$, y el espacio de búsqueda se reduce al segmento $[l,m_2]$.
	\item $f(m_1) = f(m_2)$ Podemos ver que ambos puntos pertenecen al área donde se 
	maximiza el valor de la función, o $m_1$ está en la zona de valores crecientes y 
	$m_2$ está en el área de valores descendentes (aquí usamos el rigor de la función 
	creciente/decreciente). Así, el espacio de búsqueda se reduce a $[m_1, m_2]$. Para 
	simplificar el código, este caso se puede combinar con cualquiera de los casos 
	anteriores.
\end{enumerate}

Por lo tanto, basándose en la comparación de los valores en los dos puntos internos, 
podemos reemplazar el intervalo actual $[l,r]$ con un nuevo intervalo más corto
$[l^\prime, r^\prime]$. Aplicando repetidamente el procedimiento descrito al 
intervalo, podemos obtener un intervalo arbitrariamente corto. Con el tiempo, su 
longitud será menor que una determinada constante predefinida (precisión) y el proceso 
podrá detenerse. Este es un método numérico, por lo que podemos suponer que después de 
eso la función alcanza su máximo en todos los puntos del último intervalo $[l,r]$. 
Sin pérdida de generalidad, podemos tomar $f(l)$ como valor de retorno.

No impusimos ninguna restricción en la elección de puntos $m_1$ y $m_2$. Esta elección definirá la tasa de convergencia y la precisión de la implementación. La forma más común es elegir los puntos para que divida el intervalo $[l,r]$ en tres partes iguales. Así, tenemos:

$$m_1 = l + \frac{(r - l)}{3}$$

$$m_2 = r - \frac{(r - l)}{3}$$

Si $m_1$ y $m_2$ se eligen para que estén más cerca entre sí, la tasa de convergencia aumentará ligeramente.

\subsection{El caso de los argumentos enteros}
Si $f(x)$ toma un parámetro entero, el intervalo $[l,r]$ se vuelve discreto. Dado que 
no impusimos ninguna restricción en la elección de puntos $m_1$ y $m_2$, la exactitud 
del algoritmo no se ve afectada $m_1$ y $m_2$ todavía se puede elegir dividir $[l,r]$ 
en 3 partes aproximadamente iguales.

La diferencia se produce en el criterio de parada del algoritmo. La búsqueda ternaria 
tendrá que detenerse cuando $(r-l)<3$, porque en ese caso ya no podemos seleccionar 
$m_1$ y $m_2$ ser diferentes entre sí y también de $l$ y $r$, y esto puede causar un 
bucle infinito. Una vez $(r-l)<3$, el grupo restante de puntos candidatos $(l,l+1,\ldots,r)$. Es necesario comprobarlo para encontrar el punto que produce el valor máximo $f(x)$.
\section{Implementación}
\subsection{C++}
\begin{lstlisting}[language=C++]
double f (double x){
   //Evaluacion de x en la funcion
}

double ternary_search(double l, double r) {
   double eps = 1e-9; //establezca el limite de error aqui
   while (r - l > eps) {
      double m1 = l + (r - l) / 3;
      double m2 = r - (r - l) / 3;
      double f1 = f(m1); //evalua la funcion en m1
      double f2 = f(m2); //evalua la funcion en m2
      if (f1 < f2) l = m1;
      else r = m2;
   }
   return f(l); //devolver el maximo de f(x) en [l, r]
}	
\end{lstlisting}


\subsection{Java}
\begin{lstlisting}[language=Java]
public static double f (double x){
   //Evaluacion de x en la funcion
}

public static double ternary_search(double l, double r) {
   double eps = 1e-9; //establezca el limite de error aqui
   while (r - l > eps) {
      double m1 = l + (r - l) / 3;
      double m2 = r - (r - l) / 3;
      double f1 = f(m1); //evalua la funcion en m1
      double f2 = f(m2); //evalua la funcion en m2
      if (f1 < f2) l = m1;
      else r = m2;
   }
   return f(l); //devolver el maximo de f(x) en [l, r]
}
\end{lstlisting}

\section{Aplicaciones}
A continuación, se presentan algunas aplicaciones comunes de la búsqueda ternaria:
\begin{itemize}
	\item \textbf{Optimización en funciones unimodales:} La búsqueda ternaria se utiliza para encontrar el máximo o mínimo de una función unimodal (una función que tiene un solo máximo o mínimo) en un intervalo dado. Dividiendo el intervalo en tres partes y comparando los valores de la función en los puntos intermedios, se puede encontrar el punto óptimo con mayor precisión que con la búsqueda binaria.
	\item \textbf{Búsqueda en funciones cóncavas o convexas:} En funciones cóncavas o convexas, la búsqueda ternaria puede utilizarse para encontrar el punto en el que la función alcanza su máximo o mínimo. Dividiendo el intervalo en tres partes y determinando en qué dirección se encuentra el punto óptimo, se puede reducir el espacio de búsqueda de manera eficiente.
	\item \textbf{Determinación de umbrales:} La búsqueda ternaria también se utiliza para determinar umbrales o puntos críticos en una función. Por ejemplo, si se desea encontrar el punto en el que una función supera cierto umbral, la búsqueda ternaria puede ser útil para acercarse a ese valor de manera eficiente.
	\item \textbf{Problemas de optimización numérica:} En problemas de optimización numérica donde se busca maximizar o minimizar una función en un intervalo, la búsqueda ternaria puede ser una herramienta útil para encontrar soluciones cercanas al óptimo.
\end{itemize}
 
En resumen, la búsqueda ternaria es una técnica de búsqueda eficiente que puede aplicarse en una variedad de situaciones donde se necesita encontrar el máximo o mínimo de una función, determinar umbrales o puntos críticos, o resolver problemas de optimización numérica. Su capacidad para dividir el espacio de búsqueda en tres partes puede conducir a una convergencia más rápida y precisa en comparación con otros métodos de búsqueda.
\section{Complejidad}
Se puede visualizar de la siguiente manera: cada vez que se evalúa la función en los 
puntos $m_1$ y $m_2$, esencialmente estamos ignorando aproximadamente un tercio del 
intervalo, ya sea el izquierdo o el derecho. Por tanto, el tamaño del espacio de 
búsqueda es ${2n}/{3}$ del original.

Aplicando el teorema de Maestro , obtenemos la estimación de complejidad deseada:

$$O(n) = O({2n}/{3}) + O(1) = O(\log n)$$
\section{Ejercicios}
A continuación una lista de ejercicios que se pueden resolver aplicando este algoritmo:

\begin{itemize}
	\item \href{https://codeforces.com/problemset/problem/439/D}{Codeforces - Devu and his Brother}
	\item \href{https://codeforces.com/contest/1355/problem/E}{Codeforces - Restorer Distance}
	\item \href{http://codeforces.com/problemset/problem/578/C}{Codeforces - Weakness and Poorness}
	\item \href{https://www.codechef.com/problems/AMCS03}{Codechef - Race time}
	\item \href{https://www.codechef.com/problems/ICM2003}{Codechef - Is This JEE }
	\item \href{https://acm.timus.ru/problem.aspx?space=1&num=1719}{TIMUS 1719 Kill the Shaitan-Boss}
	\item \href{https://acm.timus.ru/problem.aspx?space=1&num=1913}{TIMUS 1913 Titan Ruins: Alignment of Forces}
	\item \href{https://www.hackerearth.com/problem/algorithm/rescuer-2d2495cb/}{Hackerearth - Rescuer}
	\item \href{http://www.spoj.com/problems/KOPC12A/}{SPOJ - Building Construction}
	\item \href{http://codeforces.com/gym/101309}{GYM - Dome of Circus (D)}
	\item \href{http://codeforces.com/gym/100829}{GYM - Chasing the Cheetahs (A)}
	\item \href{https://uva.onlinejudge.org/index.php?option=com_onlinejudge&Itemid=8&page=show_problem&problem=4898}{UVA - Galactic Taxes}
	\item \href{https://uva.onlinejudge.org/index.php?option=com_onlinejudge&Itemid=8&page=show_problem&problem=3349}{UVA - 12197 - Trick or Treat}
	\item \href{http://lightoj.com/volume_showproblem.php?problem=1146}{LOJ - Closest Distance}
\end{itemize}
%\section{Anexos}
%\input{tex_files/anexos_successions.tex}


\end{document}